The blank-input reduction also combines with the root-elimination method
of packages 03 and 05. This removes one supplied unknown, at the cost
of a higher-degree equation. In this section \(i\) has its
\emph{normalized} meaning: the root to be eliminated would satisfy
\(f^2=1+\Delta i^2c^4\). The ordinary multiplier used in
Section~\ref{sec:arithmetic-proof} will be \(\Delta i\).

Delete \(f\) from the fifteen-variable list and keep all definitions
that do not involve it. Use the old positive transport coordinate
\(z_0\), so that \(N_t=(K+w)C+q-F-z_0(q-1)\). Define
\begin{align}
 Q_*&=(\Delta ic^2)^2,\qquad d_*=1+\Delta i^2c^4,\notag\\
 b_*&=c+Rd_*,\qquad v_*=cT,\notag\\
 A_*&=Q_*(d_*v_*^2+b_*^2-y^2)+y^2-1,\notag\\
 B_*&=2Q_*v_*b_*,\qquad K_*=A_*^2-d_*B_*^2,\notag\\
 P_4&=N_0N_1N_kN_t.\label{eq:rootdata14}
\end{align}
All these are polynomial abbreviations in fourteen supplied unknowns;
none introduces another existential coordinate. The first equation is
\begin{equation}\boxed{\Phi_{14}=P_4(1+K_*)-1=0.}
\label{eq:phi14}\end{equation}
Its circuit has 75 operations and exact degree 293.

A second finalizer incorporates the outer product inside the norm.
Put
\begin{align*}
 A_{\mathrm{out}}
   &=P_4\bigl[Q_*(d_*v_*^2+b_*^2)+(1-Q_*)y^2\bigr]-1,\\
 B_{\mathrm{out}}&=P_4B_*.
\end{align*}
Then
\begin{equation}\boxed{\Psi_{14}
 =A_{\mathrm{out}}^2-d_*B_{\mathrm{out}}^2=0}
\label{eq:psi14}\end{equation}
has 73 operations and exact degree 342. Both equations initially use
fourteen strictly positive unknowns. Their coefficients are the same
four fixed integers as before.

The circuits compute the norm cross-term through
\[
 d_*B_*^2=4(Q_*d_*v_*^2)(Q_*b_*^2),
\]
and its analogous outer-product expression. Thus they do not pay to
construct \(B_*\) separately. Neither circuit contains division,
square-root extraction, variable exponentiation, or a free program index.

\subsection{Restoring an integer root, with its sign left free}

At a zero of \eqref{eq:phi14}, both \(P_4\) and \(1+K_*\) are
integer units. Because \(B_*\) is even,
\[
 1+K_*\equiv1+A_*^2\pmod4,
\]
which cannot be \(-1\). Hence \(K_*=0\) and \(P_4=1\).
Every factor inside \(P_4\) is consequently a unit. The first and
main norms equal \(1\) by the descent and congruence arguments already
given. The transport-unit argument yields \(C\ge0\), and the
pre-decoding bounds \eqref{eq:Rbounds} apply; in particular \(R>0\).
On the positive domain, \(Q_*,v_*,b_*>0\), so \(B_*\ne0\).

The rational number
\[
 f=\frac{A_*}{B_*}
\]
satisfies \(f^2=d_*\). A rational square root of an integer is an
integer: write it in lowest terms, and any prime in its denominator
would have negative valuation in its square. Since \(d_*>1\),
this restored integer \(f\) is nonzero. Its sign need not be determined.
Set \(V=v_*f-b_*\). Using \(f^2=d_*\), the exact remainder identity is
\begin{equation}
 A_*-B_*f=Q_*(V^2-y^2)+y^2-1=0.
 \label{eq:root-remainder14}
\end{equation}
Moreover
\[
 V=c(Tf-1)-Rf^2,\qquad f^2-\Delta i^2c^4=1.
\]
Thus an integer solution of the normalized blank equations has been
restored, allowing a signed \(f\). Replace the multiplier by
\(i_{\mathrm{ordinary}}=\Delta i>0\). Then
\[
 (i_{\mathrm{ordinary}}c^2)^2
   =\Delta(f^2-1)=Q_*.
\]
All factors of \eqref{eq:p69} have the required product, so the
signed-domain theorem of Section~\ref{sec:arithmetic-proof} recovers
a genuine K215 blank halting run. No separate positivity theorem for
the reconstructed root is required.

For \eqref{eq:psi14}, first note that \(P_4=0\) would make its
left side equal \(1\). Otherwise reduction modulo \(P_4\) gives
\(\Psi_{14}\equiv1\pmod{P_4}\), so a zero forces
\(P_4=\pm1\). Its four factors are units and give the same
outer bounds. Since \(B_{\mathrm{out}}\ne0\), set
\(f=A_{\mathrm{out}}/B_{\mathrm{out}}\). Again \(f^2=d_*\)
makes \(f\) an integer. The corresponding remainder identity is
\[
 A_{\mathrm{out}}-B_{\mathrm{out}}f
  =P_4\bigl[Q_*(V^2-y^2)+y^2\bigr]-1=0.
\]
After the same change of strong multiplier, this is exactly the
product needed by \eqref{eq:p69}. The argument does not assume the
sign of \(P_4\) in advance.

Conversely, a canonical positive blank compiler solution has
\(f^2=d_*\), \(A_*=B_*f\), and \(P_4=1\). Deleting \(f\)
therefore makes both \eqref{eq:phi14} and \eqref{eq:psi14} vanish.
This proves the two directions of their existence equivalence and
hence their stated \(\Con(\mathrm{ZFC})\) equivalences.

\subsection{A signed extension and 32 unrestricted integer unknowns}

The nine coordinates
\[
 J,F,\alpha,z_0,i,s,w,\eta,\zeta
\]
can remain positive while the other five,
\(h,T,\tau,y,\gamma\), are arbitrary integers. The only new
obstruction to the preceding argument would be \(T=0\), making
the reconstructed-root denominator vanish. We exclude it before
taking a quotient.

For either polynomial the outer unit argument already gives
\(0<R<a<\Delta\) and \(c>2\). Since
\(Q_*=\Delta(d_*-1)\),
\begin{align*}
 Q_*-b_*&=(\Delta-R)d_*-\Delta-c\\
 &\ge d_*-\Delta-c
   =1+\Delta(i^2c^4-1)-c>0.
\end{align*}
Consequently \(2<b_*<Q_*\). If \(T=0\), then \(v_*=B_*=0\).
For \(\Phi_{14}=0\), the equation \(K_*=0\) forces \(A_*=0\),
and therefore
\begin{equation}
 Q_*b_*^2-(Q_*-1)y^2=1.\label{eq:zero-T-norm}
\end{equation}
For \(\Psi_{14}=0\), the equation
\(A_{\mathrm{out}}=0\) says that the left side of
\eqref{eq:zero-T-norm} equals \(P_4=\pm1\). It cannot equal
\(-1\), since \(Q_*\) is a square and this expression is a
square modulo four. Thus \eqref{eq:zero-T-norm} follows again.

Every solution of this equation with \(b_*>1\) has
\(b_*\ge2Q_*-1\). Indeed,
\[
 b_*^2-1=(Q_*-1)(y^2-b_*^2)
\]
implies \(|y|\ge b_*+1\), and for \(1<b_*\le2Q_*-2\)
one has \(b_*^2-1<(Q_*-1)(2b_*+1)\), a contradiction.
This conflicts with \(b_*<Q_*\). Hence \(T\ne0\), the
root denominators are nonzero, and the signed reconstruction is valid.
The signed ordinary theorem then supplies soundness exactly as above.

Use the positivity polynomial \eqref{eq:pos} independently for the
nine guarded coordinates and leave the other five unrestricted.
The resulting polynomials have \(9\cdot3+5=32\) integer unknowns.
The conversion is surjective onto the required positive values, so it
preserves existence in both directions. Their operation counts and
degrees are included below. They improve the variable count of the
explicit family studied here; they do not assert a global minimum over
all Diophantine representations.

\subsection{Exact degrees and operation ledgers}

Write \(Q_0=\beta J\), \(k_0=\eta+\zeta\), and
\(T_2=w(Q_0-F-\alpha)-z_0Q_0\). The two leading homogeneous
forms are respectively
\begin{align}
 (\Phi_{14})_{293}
  &=8h\gamma i^{12}k_0^{27}w^{21}s^{48}
       Q_0^{177}(Q_0-F)^4T_2,\label{eq:phi14-leader}\\
 (\Psi_{14})_{342}
  &=64h^2\gamma^2 i^{12}k_0^{30}w^{26}s^{56}
       Q_0^{206}(Q_0-F)^4T_2^2.\label{eq:psi14-leader}
\end{align}
They are nonzero for every \(\beta>0\). For example, the coefficients
at
\[
 J^{182}h\gamma i^{12}\eta^{27}w^{21}s^{48}z_0,
 \qquad
 J^{212}h^2\gamma^2i^{12}\eta^{30}w^{26}s^{56}z_0^2
\]
are \(-8\beta^{182}\) and \(64\beta^{212}\), respectively.

To check the degree accounting, \(Q_*,d_*,b_*,v_*\) have degrees
\(46,34,38,6\). The unique top term of \(A_*\) is \(Q_*b_*^2\),
of degree 122. Hence \(K_*\) has degree 244: the competing
term \(d_*B_*^2\) has degree 214. The degree-49 outer product gives
\(244+49=293\). In the second construction the squared leading
term has degree 342 and its subtracted cross-term degree 312.

For the integer versions assign weight two to the nine guarded
coordinates and weight one to the five signed coordinates. The leading
forms above have signed weights two and four. Their exact degrees
therefore become \(2\cdot293-2=584\) and
\(2\cdot342-4=680\). Direct weighted upper bounds give the same
values. Disjoint quadratic positivity substitutions preserve
nonvanishing, as explained in Section~\ref{sec:integer}.

\begin{center}\small
\begin{tabular}{@{}lrrrr@{}}
\toprule
Presentation & Domain & Unknowns & Degree & Operations\\
\midrule
\(\Phi_{14}\) & positive &14&293&75 = 44M+31A\\
\(\Psi_{14}\) & positive &14&342&73 = 44M+29A\\
\(\Phi_{14}\) after positivity substitution & integers &32&584&138\\
\(\Psi_{14}\) after positivity substitution & integers &32&680&136\\
\bottomrule
\end{tabular}
\end{center}
The integer conversion adds seven gates for each of nine positivity
substitutions. The files \code{arithmetic/build_fourteen.py} and its
emitted packets give complete evaluators and reproduce the counts,
root-remainder identities, substitution identities, and degree
diagnostics. These finite checks supplement the proofs; they do not
replace the unbounded compiler or Pell arguments.
